%=====================================================================
\section{Core Algorithms}
\label{app:code}
%=====================================================================

This appendix provides pseudocode for the key algorithms implementing DRLT. The full implementation is available in \texttt{lib/drlt.py}.

\subsection{Fundamental operations}

\begin{verbatim}
VERTEX(psi):
    Store psi in C^5, normalize: psi <- psi / ||psi||

OVERLAP(v_i, v_j):
    return sum(conj(psi_i[a]) * psi_j[a], a=0..4)   # G_ij in C

W(v_i, v_j):
    return |OVERLAP(v_i, v_j)|^2 / 5                 # real shadow

PHASE(v_i, v_j):
    return arg(OVERLAP(v_i, v_j))                     # gauge connection

DS2(v_i, v_j):
    return 1 - 5 * W(v_i, v_j)                       # metric
\end{verbatim}

\subsection{Network and Gram matrix}

\begin{verbatim}
NETWORK(N):
    vertices <- [VERTEX(random_C5()) for i in 1..N]

G_MATRIX(net):
    G[i,j] <- OVERLAP(vertices[i], vertices[j])      # N x N complex
    # Properties: Hermitian, PSD, rank <= 5, unit diagonal

W_MATRIX(net):
    return |G|^2 / 5                                  # N x N real

HOLONOMY(net, i, j, k):
    return arg(G[i,j] * G[j,k] * G[k,i])            # gauge flux
\end{verbatim}

\subsection{Four forces from one inner product}

\begin{verbatim}
INTERACTION_DECOMPOSITION(v_i, v_j):
    o_full <- OVERLAP(v_i, v_j)                       # full C^5
    o_S <- sum(conj(psi_i[a]) * psi_j[a], a=0..2)    # spatial C^3
    o_T <- sum(conj(psi_i[a]) * psi_j[a], a=3..4)    # temporal C^2

    gravity  <- |o_full|^2 / 5      # all d components
    strong   <- |o_S|^2 / 3         # SU(3) sector
    weak     <- |o_T|^2 / 2         # SU(2) sector
    em_phase <- arg(o_T) - arg(o_S) # U(1) relative phase

    return {gravity, strong, weak, em_phase}
\end{verbatim}

\subsection{Local evolution (tick)}

\begin{verbatim}
TICK(net):
    for each vertex i:
        # Local Hamiltonian
        H_i <- sum(W[i,j] * |psi_j><psi_j|, j != i)  # 5x5 Hermitian

        # Local effective Planck constant
        hbar_i <- LOCAL_HBAR_EFF(net, i)

        # Unitary evolution (one information bit)
        U_i <- matrix_exp(-i * H_i / hbar_i)
        psi_i <- U_i * psi_i
        psi_i <- psi_i / ||psi_i||                     # re-normalize

LOCAL_HBAR_EFF(net, i):
    A <- mean(DS2(v_i, v_j) for j != i)               # local area
    S <- binary_entropy(W[i,:])                        # local entropy
    return A / (4 * S)                                 # area per bit
\end{verbatim}

\subsection{Pachner moves}

\begin{verbatim}
PACHNER_1TO5(net):
    # Add vertex: resolution increase
    if network_diversity(net) > threshold:
        psi_new <- random perturbation of random existing vertex
        add VERTEX(psi_new) to net

PACHNER_5TO1(net):
    # Merge vertices: resolution decrease
    find pair (i,j) with W[i,j] > merge_threshold
    if found:
        psi_merged <- (psi_i + psi_j) / ||psi_i + psi_j||
        replace (i,j) with single VERTEX(psi_merged)
\end{verbatim}

\subsection{Simplex discovery}

\begin{verbatim}
FIND_SIMPLICES(net, w_threshold=0.1):
    simplices <- []
    for each 5-subset {a,b,c,d,e} of vertices:
        if min(W[i,j] for all pairs in subset) > w_threshold:
            simplices.append({a,b,c,d,e})   # emergent 4-simplex
    return simplices
\end{verbatim}

\begin{remark}
The entire physical content of DRLT is encoded in \texttt{OVERLAP}: a single inner product in $\CC^5$. Every other operation --- metric, forces, evolution, topology change --- is derived from this one function.
\end{remark}
